\documentclass[11pt]{article}
\usepackage{amsmath,amssymb,amsthm,mathtools}
\usepackage[hidelinks]{hyperref}
\usepackage[margin=1.1in]{geometry}
\newtheorem{theorem}{Theorem}
\newtheorem{lemma}{Lemma}
\providecommand{\tightlist}{\setlength{\itemsep}{0pt}\setlength{\parskip}{0pt}}
\renewcommand{\labelitemi}{$\bullet$}
\DeclareMathOperator{\Eff}{Eff}
\DeclareMathOperator{\Mov}{Mov}
\DeclareMathOperator{\Nef}{Nef}
\DeclareMathOperator{\Pic}{Pic}
\DeclareMathOperator{\Alb}{Alb}
\DeclareMathOperator{\Int}{Int}
\DeclareMathOperator{\PsAut}{PsAut}
\title{The Strong Relative Movable Cone Conjecture in Relative Dimension at Most Two}
\author{}
\date{}
\begin{document}
\maketitle
\begin{abstract}
We prove the strong relative Morrison--Kawamata movable cone conjecture for terminal
Calabi--Yau fibre spaces of relative dimension at most two. The total space may have
arbitrary dimension, the base need not be $\mathbb Q$-factorial, and the closure of the
relative movable cone may contain lines. A rational polyhedral chamber in the actual
effective movable cone covers it under support-preserving pseudo-automorphisms;
every nonidentity closed translate misses its interior and its stabilizer is trivial.
The proof combines relative log model geography with descent of a surface chamber
to the generic field. A positive-semidefinite operator construction then gives a
finite polyhedral description on every effective boundary face.
\end{abstract}

\section*{Introduction}
We establish the strong relative movable cone conjecture for terminal Calabi--Yau
fibre spaces of relative dimension at most two. The result treats arbitrary total
dimension and arbitrary admissible base and includes the degenerate case in which
the closed relative movable cone contains lines. The fundamental domain lies in
the actual effective cone, and its translates satisfy separation from the whole
closed chamber. These assertions sharpen the weak relative domains obtained in
the preceding work of Li and Zhao~\cite{LiZhao} and Li~\cite{LiII}.

The distinction between an effective cone and its rational hull is essential in
the presence of lineality. Let $C=\overline{\Mov}(X/S)\cap\Eff(X/S)$ and let
$\Gamma$ be the faithful image of support-preserving pseudo-automorphisms. We
construct a cone $\Pi\subset C$ for which $C=\Gamma\Pi$ and
$\gamma\Pi\cap\Int(\Pi)=\varnothing$ whenever $\gamma\ne1$. Thus the domain
retains the effective boundary and has no residual finite stabilizer.

Three ingredients make the passage from weak coverage to this strong statement.
First, rational approximation of the boundary and relative model geography
control effective movable classes and the vertical directions. Second, descent
from the complex surface theorem produces a chamber for the automorphisms
actually defined over the generic field, while the log pluricanonical action
identifies a finite-index subgroup extending over the base. Third, the Picard
quotient over a numerical class $d$ has the exact affine fibre
\[
  s(d)+\operatorname{im}T_d,
\]
where $T_d$ is a rational positive-semidefinite operator. A finite system of
inequalities cuts out the chamber and removes directions in $\ker T_d$ on each
boundary face. The argument uses the surface cone theorem of Totaro~\cite{Totaro}
and the geometry of polyhedral-type group actions developed by
Looijenga~\cite{Looijenga}.

Here is the precise result. We use the conventions of Li and Zhao for a terminal
Calabi--Yau fibre space and for the support-preserving pseudo-automorphism group.

\begin{theorem}\label{thm:main}
Let $f\colon(X,\Delta)\to S$ be a complex terminal Calabi--Yau fibre space in the sense of Li and Zhao, where $X$ is normal, projective and $\mathbb Q$-factorial, $K_X+\Delta\equiv_S0$, and $0\leq\dim X-\dim S\leq2$. Set
\[
 V=N^1(X/S)_{\mathbb R},\qquad
 C=\overline{\Mov}(X/S)\cap\Eff(X/S),\qquad
 \Gamma=\operatorname{Im}\bigl(\PsAut(X/S,\Delta)\longrightarrow\operatorname{GL}(V)\bigr).
\]
There is a rational polyhedral cone $\Pi\subset C$ such that $\Gamma\Pi=C$ and
\[
 \gamma\Pi\cap\Int(\Pi)=\varnothing\qquad(\gamma\in\Gamma\setminus\{1\}).
\]
The setwise stabilizer of $\Pi$ in $\Gamma$ is trivial. No assumption of nondegeneracy of $\overline{\Mov}(X/S)$ or of $\mathbb Q$-factoriality of $S$ is imposed.
\end{theorem}
\begin{proof}
\subsection*{1. Relative models and vertical classes}

Write \(V\) for the real Cartier numerical group for all complete curves contracted by \(f\), with its
usual rational structure. Write \(M\) for the cone generated by relatively movable Cartier divisors,
\(E\) for the cone of relatively effective divisors, and \(C\) for the intersection of the closure of
\(M\) with \(E\). Write \(\Gamma\) for the faithful linear image of the pseudo-automorphisms over \(S\)
preserving the support of \(\Delta\). Precomposition with the inversion map permits computations with
pullback instead of pushforward; all sets of linear operators or translates below are the same. The base
\(S\) is normal quasi-projective. With \(X\) projective and \(f\) surjective it is proper and hence
projective. No divisor on \(S\) is assumed \(\mathbb Q\)-Cartier unless specifically chosen to be a
Cartier divisor.

For a terminal fibre space it is the pair \((X,\Delta)\) which is terminal in the terminology of Li--Zhao
Section~2.2; see also Li II Section~2.2, which specifies \(\Delta\) effective and log discrepancy
$>$1 at all exceptional divisorial valuations. In dimension at least two the coefficients of
\(\Delta\) are $<$1: at the general smooth locus of any chosen component blow up a
codimension-two smooth centre contained in it. This gives log discrepancy 2 minus that coefficient. All
remaining log discrepancies are positive, so the pair is klt. Absolute dimension zero gives \(V=0\). In
absolute dimension one, a morphism to a point has the single nonnegative degree ray with trivial linear
image, and a nonconstant contraction of normal curves is an isomorphism. These cases also cover any
low-dimensional convention about terminal boundary coefficients. From now on we use klt.

Fix the finitely many prime divisors supporting \(\Delta\). Relative numerical triviality of \(K_X\) plus
a boundary supported on them is a rational affine system in its coefficients. Indeed all these divisors
are \(\mathbb Q\)-Cartier, and the image of their span in the finite-dimensional numerical space is
rational. Terminality is open under small perturbations on this fixed support (the discrepancy criterion
on a fixed log resolution, Kollár--Mori Cor. 2.35). Choose rational points of this affine system in a
small neighborhood whose convex hull contains \(\Delta\), all positive on its support and terminal. For
each resulting boundary \(\Delta'\) the rational relative numerical-triviality theorem of
Hacon--Xu~\cite{HaconXuNumerical} gives \(K_X+\Delta'\) relatively \(\mathbb Q\)-linearly trivial (Cor.
1.4 in Li--Zhao's published numbering, Cor. 1.6 of arXiv:1107.4149v2). Taking convex combinations gives
relative \(\mathbb R\)-linear triviality of \(K_X+\Delta\). Fix one such rational boundary \(\Delta_0\)
with \(\operatorname{Supp}\Delta_0=\operatorname{Supp}\Delta\).

For all subsequent geography, good-model and group arguments use the working pair \((X,\Delta_0)\). The
original \(V\), \(M\), \(E\), and \(C\) depend only on \(X/S\). The group of relative
pseudo-automorphisms preserving the support is literally the same for \(\Delta_0\) and the given real
\(\Delta\), and its linear action on \(V\) is the same. Consequently a domain constructed for this
working pair answers the theorem for the original real pair. A component-fixing finite-index subgroup,
when needed, preserves the coefficients of \(\Delta_0\) as well. The effective directions used to perturb
\(\Delta_0\) can still have arbitrary real coefficients.

The relative good-model and geography results of Li--Zhao apply in relative dimension at most two. On
very general closed surface fibres Fujino~\cite{Fujino}, \emph{Minimal model theory for log surfaces},
Thm. 1.1, gives good minimal models for effective klt pairs, including real-boundary perturbations; the
curve and point cases are elementary. Li--Zhao~\cite{LiZhao}, Thm. 2.4 and its explanation after the
statement give the klt real-coefficient relative criterion. Hacon--Xu~\cite{HaconXuGood}, Thm. 2.12, from
which it is derived, uses a rational boundary and finite generation. In the real extension the needed
relative log canonical model is provided by Z. Li~\cite{LiFiniteness}, \emph{On finiteness of log
canonical models}, Cor. 1.2, under the hypothesis of ordinary relative Kodaira dimension at
least zero. That hypothesis holds here by an integral round-down computation.

Indeed there exist a positive integer \(r\), \(\varphi\in\mathbb C(X)^*\), and an integral Cartier divisor \(B_S\) on \(S\) such that \(r(K_X+\Delta_0)=\operatorname{div}_X(\varphi)+f^*B_S\). Fix any effective real Weil divisor \(D\) and any real \(\epsilon\ge0\) for which \((X,\Delta_0+\epsilon D)\) is klt. For every \(m=r\ell\), \(\ell\) a positive integer, equality of Weil divisors gives \[
 \left\lfloor m(K_X+\Delta_0+\epsilon D)\right\rfloor
 =\ell\operatorname{div}_X(\varphi)+\ell f^*B_S+
   \left\lfloor m\epsilon D\right\rfloor .
\] The final summand is effective. Over an open set on which \(B_S=\operatorname{div}(h)\), the rational function \((\varphi f^*h)^{-\ell}\) is a nonzero section of the sheaf of this rounded divisor. Hence \(f_*\mathcal O_X(\lfloor m(K_X+\Delta_0+\epsilon D)\rfloor)\) is nonzero at the generic point for every such \(m\). This proves \(\kappa(K_X+\Delta_0+\epsilon D;X/S)\ge0\) by Li's Definition 3.1, also when \(D=0\) or \(D\) is vertical. The adjoint on a very general closed fibre is pseudo-effective and has a good model by the relative-dimensional result just cited. The two inputs to the relative argument are thus available for each actual effective real perturbation. The following consequences are needed.



\emph{Effective movability.} We claim that \(M=C\). For an effective representative \(D\) of a class in
the closed movable cone, take the good model of \((X,\Delta_0+\epsilon D)\) over \(S\),
\(\epsilon>0\) small. Its existence follows from generic low-dimensional abundance and
Li--Zhao~\cite{LiZhao}, Thm. 2.4; this includes vertical \(D\) with generically zero restriction. Lem.
2.3 allows a weak log canonical model \(X'\) which is a small \(\mathbb Q\)-factorial modification. Write
\((\Delta_0)_{X'}\) and \(D_{X'}\) for the strict transforms. The corresponding log adjoint is relatively
\(\mathbb R\)-semiample on \(X'\), since the pullback by the morphism supplied by Lem. 2.3 of the adjoint
on the original good model is semiample. Strict transformation of the working relative \(\mathbb
Q\)-principal relation gives \(K_{X'}\)+\((\Delta_0)_{X'}\) relatively \(\mathbb Q\)-linearly trivial.
Thus \(D_{X'}\) is relatively \(\mathbb R\)-semiample.

A relatively \(\mathbb R\)-semiample divisor is a positive real combination, up to relative real linear
equivalence, of semiample rational classes; their transforms are relatively movable. Thus \([D]\) belongs
to \(M\). The reverse inclusion follows at once from the definitions. This is equality with the effective
cone intersection only; no equality with the rational hull is used.

\emph{Polyhedral intersections.} For each rational polyhedral cone \(P\) contained in \(E\), the
intersection \(P\cap C\) is rational polyhedral. Here is the direct geography argument, which also proves
the required instance of Li II~\cite{LiII}, Lem. 2.5(2). Li--Zhao~\cite{LiZhao}, Thm. 2.7 decomposes a
pointed \(P\) into finitely many relatively open rational polyhedral cells on which effective small klt
perturbations share their marked weak log canonical models. In any such cell meeting \(C\) take one
\([D]\in C\) and the small marked model for \(D\) from the preceding paragraph. For a small
marking the unperturbed adjoints have the same pullback on a common resolution by their common relative
\(\mathbb Q\)-principal relation with compatible canonical divisors.

The discrepancy differences for the perturbation of \(D\) therefore scale linearly with
\(\epsilon>0\), so the marked model stays weak log canonical when \(\epsilon\) decreases.
For every effective \(B\) in the same cell, take a common perturbation small enough for \(D\) and \(B\).
The same marking makes \(B\) nef, since its transformed log Calabi--Yau summand is relatively \(\mathbb
R\)-linearly trivial. A relatively nef class on the marked model is a limit of relatively ample classes,
and their transforms on \(X\) are movable. Hence \([B]\) is in the closed movable cone and, as
\([B]\in P\subset E\), it belongs to \(C\). The closure of this entire cell still lies in
the closed movable cone and in \(P\subset E\).

Thus \(P\cap C\) is precisely the finite union of the rational polyhedral closures of the cells meeting
it. Convexity makes this union their Minkowski sum. For a \(P\) containing lines, intersect \(P\) with
the finitely many coordinate orthants of any rational basis of \(V\). This writes \(P\) as a finite union
of pointed rational polyhedral cones. Apply the two results to each piece, where the convex hull of a
suitably normalized finite set of nonzero ray generators avoids the origin, as required in the proof of
Li--Zhao~\cite{LiZhao}, Thm. 2.7. A zero piece needs no decomposition.

The finite union of all the resulting \(P_i\cap C\) equals the convex \(P\cap C\), hence equals
their Minkowski sum and is rational polyhedral. Finiteness of the marked models used below follows by
taking the union of the finite lists of cells on the pointed pieces.



Let \(T\subset V\) be the real span of the classes of all vertical prime divisors, and \(p\colon V\to
L=V/T\). This is a rational subspace. For a vertical prime \(D\), an effective Cartier divisor on \(S\)
containing its image has pullback containing \(D\). Thus a multiple of that pullback minus \(D\) is
effective. In particular \(T\subset E\), with both signs. The action on \(T\) has finite image: choose
finitely many vertical primes whose classes span it. The transforms of any one of these primes are among
the finitely many divisorial components of the fixed closed subset \(f^{-1}(f(D))\), since \(f\) is
proper and all maps act over the identity of \(S\).

\textbf{Finite vertical kernel.} The subgroup \(K=\ker(\Gamma\to\operatorname{GL}(L))\) is finite. Pass
to a finite-index subgroup \(K_v\) fixing \(T\) pointwise that is induced by pseudo-automorphisms fixing
each component of the common support, hence its working coefficient in \(\Delta_0\). This passage uses
only finiteness of that support. If \(A\) is an effective relatively ample Cartier class, every \(gA\) for
\(g\in K_v\) belongs to the rational polyhedral cone \(Q=\mathbb R_{\ge0}A+T\subset E\). The
marked nef cones \(g\operatorname{Nef}(X/S)\) occurring here are finite in number. Indeed, use the finite geography on
\(Q\) above. A marked small modification taking gA to an ample divisor is a weak log canonical model of
the pair perturbed in the effective direction gA.

To see this, on a common resolution the working log divisor is crepant, since strict transformation
preserves \(\Delta_0\) and its relative \(\mathbb Q\)-principal relation. The difference between
pullbacks of the perturbing divisor and of its relatively ample transform is exceptional for both
projections and is nonnegative by the negativity lemma, applied over the original model. If gA and hA
belong to the same cell, choose a common positive perturbation small enough for both. The same marked
model then takes \(hA\) to a nef class. Thus \(hA\), an interior point of \(h\operatorname{Nef}(X/S)\), belongs to
\(g\operatorname{Nef}(X/S)\). Both cones are full dimensional, and distinct marked ample chambers have disjoint
interiors (reconstruction by relative Proj from an ample divisor, as in Kawamata~\cite{Kawamata}, Lem.
1.5).

Taking interior points of the latter cone near hA shows that the two cones coincide. A relative numerical
ample class has an effective representative because the base is quasi-projective; we chose A in this way
from the start.

Consequently the subgroup \(K_1\) of \(K_v\) stabilizing \(\operatorname{Nef}(X/S)\) has finite index. For
\(g\in K_1\), \((g-I)^2=0\). Every relatively ample class A' satisfies \(g^nA'=A'+n(g-I)A'\) in
the closed relative nef cone for all integers \(n\).~Its closure thus contains both signs of (\(g\)-I)A'.
The nef cone is pointed by the very definition of numerical equivalence, so (\(g\)-I)A'=0 for every \(A'\).
Ample classes span \(V\). This proves \(K_1=\{I\}\) and the lemma.

\textbf{Lifting a quotient chamber.} Suppose the faithful image \(G\) of \(\Gamma\) on \(L\) has a strong
rational polyhedral fundamental domain \(P\) on \(p(C)\). Every rational generator of \(P\) has a
rational effective lift. Indeed, choose any effective real lift in \(C\), express it in finitely many
integral effective divisor classes, and solve for nonnegative rational coefficients giving the specified
rational image under \(p\).~A nonempty rational polyhedron has a rational point. It follows that
\(p^{-1}(P)\) is generated by such effective lifts together with \(T\), hence is rational polyhedral and
lies in \(E\). Therefore \(D=C\cap p^{-1}(P)\) is rational polyhedral. Its \(\Gamma\)-translates cover
\(C\).

It is full dimensional since \(\Gamma\) is countable and \(C\) contains the relatively ample open cone.
Surjectivity gives \(p(D)=P\) and openness of \(p\) gives \(p(\operatorname{Int}D)\subset\operatorname{Int}P\).
By the strong separation for \(P\), a translate of \(D\) meeting Int \(D\) must project to the identity
in \(G\). The whole setwise stabilizer is exactly \(K\), finite by the preceding lemma. Refine any
full-dimensional weak cone \(D\) with finite stabilizer \(K\) by choosing a rational linear form \(\ell\) on \(V\)
whose stabilizer in \(K\) is trivial, and imposing \(\ell(x)\ge\ell(kx)\) for every \(k\in K\). Its
\(K\)-translates cover \(D\) since the maximum over \(K\) is attained; countability/full dimension
implies the refinement has interior.

On the interior every inequality indexed by \(k\ne1\) is strict because its functional is nonzero.
Comparing that inequality with the one for the inverse proves kD' avoids Int \(D\)' when \(k\ne1\).
Elements outside \(K\) already avoid Int \(D\). This proves Definition 3.4 including the trivial setwise
stabilizer.

If the relative dimension is zero, \(f\) is birational (a generically finite normal contraction in
characteristic zero is generically degree one), \(G\) and \(\Gamma\) are trivial, and \(T=V\) since every
divisor is vertical. The preceding intersection result with \(P=V\subset E\) gives \(C\) rational
polyhedral. This also deals with non-small crepant morphisms. We assume relative dimension 1 or 2 from
now on.

\subsection*{2. The generic fibre and its numerical chamber}

Let \(k=\mathbb C(S)\), \(F=X_\eta\), \(R=N^1(F)_{\mathbb Q}\), and write rational spaces below before
extending scalars. Since the working boundary \(\Delta_0\) is effective, removing it can only increase
exceptional discrepancies, so \(X\) itself is terminal. Clemens--Kollár--Mori~\cite{CKM}, \emph{Higher
dimensional complex geometry}, Cor. 6.13, gives codimension at least three for its singular locus; no
component of that locus dominates \(S\) when the relative dimension is at most two. Generic smoothness of
the smooth open part gives \(F\) geometrically smooth. Normality and geometric connectedness of the
generic fibre also follow from Li--Zhao~\cite{LiZhao}, Prop. 4.1. Restriction yields a geometrically klt
Calabi--Yau pair on \(F\). Generic pseudo-automorphisms of this smooth curve or surface are regular
automorphisms, since a birational map of smooth curves or smooth surfaces that contracts and extracts no
prime divisor is an isomorphism. Restricting a relative pseudo-automorphism and its inverse preserves the
codimension-one isomorphism at the generic fibre. Conversely we do not initially need every generic
automorphism to extend.

Restriction of divisors gives a surjection \(\operatorname{Pic}(X)_{\mathbb Q}\to\operatorname{Pic}(F)_{\mathbb Q}\) by horizontal closure and \(\mathbb Q\)-factoriality. Its kernel consists of classes with a vertical representative, because a rational function on \(F\) is one on \(X\). Let \(K_0\) denote the subspace of \(\operatorname{Pic}(F)_{\mathbb Q}\) coming from \(f\)-numerically trivial rational divisors. Such classes are numerically trivial on \(F\) by Li--Zhao~\cite{LiZhao}, Prop. 4.3 and lie in \(\operatorname{Pic}^0(F)(k)_{\mathbb Q}\): geometric numerical and algebraic equivalence of line bundles agree modulo torsion. Points of the Picard scheme versus actual line bundles defined over \(k\) differ by a torsion Brauer obstruction; allowing rational coefficients removes it. Thus \[
 L_{\mathbb Q}=\operatorname{Pic}(F)_{\mathbb Q}/K_0,\qquad
 U_{\mathbb Q}=\operatorname{Pic}^0(F)(k)_{\mathbb Q}/K_0,\qquad
 0\longrightarrow U_{\mathbb Q}\longrightarrow L_{\mathbb Q}
 \stackrel{\pi}{\longrightarrow}R\longrightarrow0.
\] In particular \(U\) is finite dimensional independently of the size of rational points of the abelian variety \(\operatorname{Pic}^0(F)\). Here and below a real Pic class denotes an element after tensoring over \(Q\) with \(R\).

\textbf{The actual projected cone.} Put
\(D_{\mathrm{gen}}=\operatorname{Nef}(F)\cap\operatorname{Eff}(F)\), where \(\operatorname{Nef}(F)\) denotes the closed
nef cone. The cone \(p(C)\) consists precisely of the images in \(L\) of effective real Pic classes on
\(F\) with nef numerical class. To prove the forward direction, restriction of a relatively movable
linear system has no fixed divisor on \(F\); it is nef on a surface because for each given curve one may
avoid it in a representative. A relatively effective Cartier divisor means one relatively \(\mathbb
R\)-linearly equivalent to an effective divisor, so its numerical representative restricts to a generic
effective class, and nefness persists under real limits.

For the converse, on a surface express an effective representative as a nonnegative real combination of
finitely many prime curves. Nefness on this support is given by the finitely many rational intersection
inequalities against the support curves (or their finitely many geometric components) and nonnegativity
of the coefficients. Every other geometric curve has nonnegative intersection with each component. Thus
it is a nonnegative sum of rational effective nef divisors on the same support by rational polyhedrality,
imposing Galois-invariant coefficient relations if using geometric components. On curves this is
immediate. Each rational summand is semiample by low-dimensional klt abundance applied to a small
effective perturbation of a rational log Calabi--Yau boundary with the same support.

The rational boundary \(\Delta_0\) of Section 1 restricts to such a boundary, so rational
\(\epsilon>0\) gives \(K_F+(\Delta_0)_F+\epsilon D\) rationally linearly equivalent to
\(\epsilon D\). Surface abundance applies after geometric scalar extension and, on curves, the
corresponding statement follows from degree or the effective degree-zero case. Semiampleness of a
multiple descends to \(k\) by base change of global sections.

For a semiample generic Cartier class choose an effective multiple \(H\) and close its horizontal support
to an effective \(\mathbb Q\)-Cartier divisor \(D\) on \(X\), with the original multiplicities. \(D\) has
no vertical component. To avoid a putative horizontal fixed prime, choose \(H\)' in the same generic
linear system as \(H\) avoiding that prime, with horizontal closure \(D\)'. Then
\(D\)-\(D\)'-div(\(\alpha\)) is vertical for a rational function \(\alpha\). An effective Cartier divisor
on \(S\) whose pullback dominates the negative part makes \(D\) relatively linearly equivalent to \(D\)'
plus an effective vertical divisor. After a Cartier multiple this yields a section in the relative system
avoiding that prime. The original \(D\) avoids every vertical prime. Hence \(D\) is movable over \(S\).
This is Li II proof of Thm. 1.3, Step 1 and works for arbitrary base and total dimension. Multiples and
real sums yield the assertion. In particular \(\pi p(C)\)=\(D_{\mathrm{gen}}\).

\textbf{Finite-index generic lifting.} Let \(J\) be the group of \(k\)-automorphisms of \(F\) preserving the common support of \((\Delta_0)_F\) and \(\Delta_F\). Take the subgroup \(J\)' fixing all of its finitely many geometric prime components. This subgroup has finite index in \(J\); it also fixes the coefficients of \((\Delta_0)_F\). Clearing denominators in the relative \(\mathbb Q\)-linear triviality from Section 1 gives an integer \(r>0\), a Cartier divisor \(D_S\) on \(S\) and a rational \(r\)-pluricanonical form \(\Phi\) on \(X\) satisfying \[
 \operatorname{div}_X(\Phi)+r\Delta_0=f^*D_S .
\] If relative triviality is defined locally on \(S\), choose a common multiple over a finite affine cover. The nowhere vanishing local generators then glue to the pullback of a line bundle on \(S\) because \(f_*\mathcal O_X=\mathcal O_S\). No Weil divisor of \(S\) is used as a Cartier divisor without this construction.

After dividing by the rth power of a fixed rational top form of \(k/\mathbb C\), \(\Phi\) determines the
generator of the one-dimensional space \(H^0(F,\omega_F^r(r(\Delta_0)_F))\). Every \(g\in J\)' is
a regular automorphism of the klt pair (\(F\),\((\Delta_0)_F\)), hence is log crepant and acts on this
space by a scalar \(\lambda(g)\in k^*\). Apply Fujino--Gongyo~\cite{FujinoGongyo}, \emph{Log
pluricanonical representations and the abundance conjecture}, Thm. 1.1 (also their klt Cor. 3.10) to the
geometric pair: it is projective klt with \(r(K_F+(\Delta_0)_F)\) Cartier and trivial. One may
identify the algebraic closure of \(k=\mathbb C(S)\) abstractly with \(C\), since both have
characteristic zero and the same transcendence degree over \(Q\), and the space of sections commutes with
this scalar extension. Its geometric log pluricanonical representation has finite image, so
\(\lambda(g)\) is a root of unity. All roots of unity in \(k\) lie in its original constant field
\(C\).

Since \(g\) fixes \(k\), passage back to the absolute form gives \(g^*\Phi=\lambda(g)\Phi\). Every divisorial valuation \(v\) of \(\mathbb C(X)\) over \(X\) is trivial on \(\lambda(g)\). Its discrepancy is \[
 a(v;X,\Delta_0)=1+\big(\operatorname{ord}_v\Phi-
                    \operatorname{ord}_v f^*D_S\big)/r ,
\] computed by pulling back to a smooth model carrying \(v\). The second order in this formula is invariant because the induced birational map is over \(S\). Therefore \(g\) preserves all discrepancies. For the effective terminal pair (\(X\),\(\Delta_0\)), the valuations with \(a\le1\) are precisely the prime divisors on \(X\). It follows that \(g\) and its inverse contract no divisor and extend to a pseudo-automorphism of \(X/S\). On these prime divisors discrepancy recovers the coefficient of \(\Delta_0\) as 1-a, so the extension preserves the support of the original \(\Delta\), including its vertical components.

This proves that the actual generic subgroup contains \(J\)'. In particular, every \(k\)-point in the
neutral support-preserving automorphism group extends, since a geometrically connected group fixes each
geometric boundary component. Only a finite-index assertion for \(J\) is used.

\textbf{A strong base chamber for the actual group.} In dimension one the generic numerical cone is the
invariant degree ray. In dimension two we give an explicit descent from the complex absolute surface
theorem to the genuine \(k=\mathbb C(S)\)-automorphisms. Write \(\Theta=(\Delta_0)_F\), let
\(W=N^1(F_{\bar k})_\mathbb R\) with its Néron--Severi lattice modulo torsion, and let \(O\subset W\) be
its open ample cone. This cone is nonempty and full dimensional, and its closure is pointed. Fix an
abstract field isomorphism \(\bar k\simeq\mathbb C\); the algebraically closed fields have the same
characteristic and transcendence degree over \(\mathbb Q\). Totaro~\cite{Totaro}, Theorem 4.1 gives a
rational polyhedral cone \(Q\) in the geometric effective nef cone whose translates under
\(M=\operatorname{Aut}(F_{\bar k},\Theta_{\bar k})\) cover that cone. Here \(M\) preserves the boundary
as a divisor. The rational generators of \(Q\) are nef, so \(Q\subset O^+\). The faithful image
\(\Gamma_{\mathrm{geom}}\) of \(M\) in \(\mathrm{GL}(W)\) preserves the lattice; it is of polyhedral type
and hence finitely generated by Looijenga, Proposition--Definition 4.1 and Corollary 4.15.

Let \(P\) be the Albanese torsor of \(F\) under \(D=\operatorname{Alb}^0(F)\), with its natural Albanese morphism. The identity component \(A=M^0\) maps isogenously onto the translations of \(P_{\bar k}\) by Jinsong Xu, Theorem 4.5. Fix a degree-\(n\) closed point of \(F\), and let \(z_1,\ldots,z_n\) be the images on \(P\) of its geometric conjugates, counting repetitions. The set \[
 Z=\left\{z\in P(\bar k):\sum_{i=1}^n(z_i-z)=0\text{ in }D(\bar k)\right\}
\] is Galois stable and is one nonempty finite coset of \(D[n]\), as follows by choosing a geometric origin and using surjectivity of multiplication by \(n\). For \(D=0\) it is the unique point. Every \(g\in M\) induces an affine map on \(P_{\bar k}\) whose linear part preserves \(D[n]\). Hence \(gZ\) is also a coset, and some element of \(A\) translates it back to \(Z\). Since connected automorphisms act trivially on the Néron--Severi group, the setwise stabilizer \(M_Z\) still surjects onto \(\Gamma_{\mathrm{geom}}\).

The kernel \(N\) of this surjection is finite. The geometric automorphisms preserving \(\Theta_{\bar k}\)
and the numerical class of a fixed integral ample line bundle form an algebraic group of finite type:
their graphs have one Hilbert polynomial for the product polarization, and the boundary conditions are
closed in the automorphism scheme. Its identity component is \(A\). Within any of its finitely many
components two elements of \(N\) differ by an element of \(A\) preserving \(Z\). That element maps under
the isogeny \(A\to D_{\bar k}\) to the finite group \(D[n]\), so there are finitely many possibilities.
Thus \(M_Z\) is finitely generated by lifts of generators of \(\Gamma_{\mathrm{geom}}\) together with
\(N\).

The absolute Galois group acts on the group \(M_Z\). Choose a finite Galois extension of \(k\) over which
its finitely many generators are defined as geometric maps and over which the geometric Néron--Severi
action is trivial. Write \(H\) for the finite Galois group of this extension over \(k\). Both actions
factor through \(H\). The fixed subgroup \(M_Z^H\) consists of actual \(k\)-automorphisms of
\((F,\Theta)\), by descent of the graphs. The semidirect product \(M_Z\rtimes H\) acts linearly on \(W\),
preserving \(O\) and its lattice. Its numerical representation has finite kernel: for each fixed \(h\in
H\), any fibre among the elements \((g,h)\) has size at most \(|N|\).

For this semidirect product the following set is finite: \[
  \mathcal S=\{a\in M_Z\rtimes H:\ Q\cap aQ\cap O\ne\varnothing\}.
\] Apply Looijenga~\cite{Looijenga}, Siegel property (Theorem 3.8) to the faithful numerical image and the finitely many pairs of faces of \(Q\) whose relative interiors meet \(O\). Such a relative interior is contained in \(O\), so its smallest containing face of \(O^+\) is \(O^+\) and the corresponding pointwise stabilizer in the faithful linear group is trivial. This gives finiteness in that group, as in Looijenga's proof of Proposition 4.6. The representation of the actual semidirect product has finite kernel, giving the displayed finiteness before passing to numerical images.

Take \(x\in O\cap W^H\). Choose \(g\in M_Z\) with \(y=g^{-1}x\in Q\). Each \(g^{-1}hg\), for \(h\in H\),
fixes \(y\in O\cap Q\), so belongs to \(\mathcal S\). Only finitely many labelled tuples
\((g^{-1}hg)_{h\in H}\) occur. Choose one \(g_j\) for each. If the tuples for \(g\) and \(g_j\) agree,
then \(g g_j^{-1}\) centralizes \(H\) in the actual semidirect product, hence is a member of \(M_Z^H\).
Consequently the \(M_Z^H\)-translates of the finitely many rational polyhedral cones \(g_jQ\cap W^H\)
cover \(O\cap W^H\).

Base change identifies \(R=N^1(F/k)_\mathbb R\) with \(W^H\) over the natural rational structures.
Numerical equivalence can be checked on geometric curves by descent to finite extensions and the
projection formula. Conversely, for a rational invariant class, choose a geometric rational divisor
representing it, descend it to a finite Galois extension and trace down; the resulting numerical class is
the extension degree times the original class. Ample rational classes are detected by field extension;
taking convex hulls of nearby rational ample classes identifies the full real ample open cone with
\(O\cap W^H\). Every rational generator of \(g_jQ\cap W^H\) has a \(k\)-effective nef rational
representative.

Indeed it lies in the geometric effective nef cone. Expressing a rational vector as a nonnegative real
combination of finitely many geometric prime classes gives a rational linear feasibility problem and
hence a nonnegative rational solution. Tracing this effective representative to \(k\) and dividing by its
finite extension degree preserves effectivity and numerical class. The convex cone \(Q_k\) generated by
the finite list of cones therefore lies in the actual \(D_{\rm gen}\), and its translates by the field
automorphisms cover the \(k\)-ample open cone.

These field automorphisms belong to the generic support-preserving group \(J\) above. Its finite-index
subgroup \(J'\) lifts to the original relative group. Choose finitely many representatives with
\(J=\bigcup_i J'g_i\), using the opposite group if needed for pullbacks. The rational polyhedral cone
generated by \(g_iQ_k\) lies in the convex \(D_{\rm gen}\). Its translates by the actual numerical image
\(\Gamma_R\) cover the \(k\)-ample open cone. Every point of \(D_{\rm gen}\) is a nonnegative combination
of rational effective nef classes by the fixed-support argument above, so \(D_{\rm
gen}\subset\operatorname{Nef}(F)^+\). That image \(\Gamma_R\) preserves the generic Néron--Severi
lattice. Li--Zhao~\cite{LiZhao}, Prop. 3.3 for the pointed generic cone gives the whole rational hull
\(\operatorname{Nef}(F)^+\) as the union of the translates of this cone. Every such translate lies in
\(D_{\rm gen}\), so the two cones agree. Their Lem. 3.5 now gives a strong rational polyhedral domain
\(B\subset D_{\rm gen}\) for the faithful \(\Gamma_R\), including separation from whole closed
translates.

If \(q=h^1(F,\mathcal O_F)=0\), \(U=0\). Thus \(p(C)=D_{\mathrm{gen}}\) has the required strong
chamber; Section 1 completes this case. In the remaining cases assume \(U\ne0\), as \(U=0\) is settled by
the same argument independently of \(q\). By Li II~\cite{LiII}, Thm. 4.1 and its descent in Section~4,
the neutral support-preserving automorphism group A of \(F\) is an abelian variety of dimension \(q\).
Its action on the Albanese torsor induces a \(k\)-homomorphism A\(\to\)Alb(\(F\)), since the geometric
identity component of the automorphisms of that torsor consists of translations. Over the geometric
closure this is the isogeny in Xu~\cite{Xu}, Thm. 4.5(2), hence it is an isogeny over \(k\).

The group preserving the reduced support has the same neutral component when the coefficients on that
support are changed. Its orbits have dimension \(q\) (the geometric presentation is a finite quotient of
A times the Albanese fibre). It follows \(q\le\dim F\le2\). If \(q=2\), \(F\) is an abelian surface
torsor: the smooth connected Albanese fibre is zero dimensional, or equivalently after algebraic closure
the proper connected orbit has full dimension and a finite stabilizer; after quotienting by this
stabilizer its underlying variety is an abelian surface. The full neutral automorphism group of an
abelian surface is its translation group, which acts simply transitively.

The effective horizontal boundary must be zero in this case, because its class equals -\(K_F\)=0
geometrically and an effective nonzero divisor on a projective surface has positive intersection with an
ample class. Descent makes \(F\) a torsor for the abelian translation group A/\(k\).

\subsection*{3. The translation subgroup}

Fix an integral relatively ample \(H_X\) on \(X\) with generic numerical class \(H\). Section 2 gives the
extension of every \(b\in\)\(A(k)\), including preservation of vertical boundary. For this neutral
subgroup its action on the one-dimensional generic log pluricanonical space can also be seen directly to
be trivial, since an abelian variety has no nontrivial algebraic character to the multiplicative group.

Write \(G\) for the faithful image on \(L\) of the actual \(\Gamma\) and \(G_0\) for its kernel on \(R\).
An operator \(g\) below is the pullback \(\varphi^*\) of a geometric map \(\varphi\); the multiplication
in \(G\) is composition of the operators, so \(\varphi^*\psi^*=(\psi\circ\varphi)^*\). Equivalently this
is the original pushforward representation after indexing \(\varphi\) by its inverse. Thus the operator
denoted \(t_b^*\) is the pullback of translation by \(b\) and is the original pushforward by translation
by -\(b\). The preimage of \(G_0\) in actual generic automorphisms preserves the numerical class of
\(H\). The subgroup scheme of support-preserving automorphisms preserving an ample numerical class is of
finite type (bounded graphs with fixed Hilbert polynomial, cf.~Li--Zhao Lem. 4.8), with identity
component A. Thus its \(k\)-points modulo their intersection with \(A(k)\) form a finite group. This
remains true for the actual subgroup and its image. The image of \(A(k)\) is therefore a normal
finite-index subgroup of \(G_0\).

For a line bundle \(M\) on \(F\) and \(b\in\)\(A(k)\), set \(\lambda_M(b)=t_b^*M-M\) in
\(\operatorname{Pic}^0(F)\). The square theorem for abelian varieties (and restriction to orbits when
\(q=1\)) shows that it depends only on the numerical class \(d\) of \(M\), is additive in \(d\), and is a
homomorphism in \(b\). For the ample \(H\) it is an isogeny A\(\to\operatorname{Pic}^0(F)\), either
by the usual polarization in the torsor case or by the orbit restriction computed below. Identify \(U\)
with its description through \(\lambda_H\).

The actual translation parameters \[
 \Lambda=\{[\lambda_H(b)]\in U\mid b\in A(k)\}
\] form a full lattice in \(U\). Their \(\mathbb Q\)-span is all \(U\) because \(\lambda_H\) is an isogeny, surjective after tensoring \(k\)-points with \(Q\): lift a rational point in a finite separable extension and sum its conjugates. Alternatively a quasi-inverse isogeny gives the same conclusion. They lie in a discrete lattice since a translation shifts \(H_X\) by an integral class of the \(\Gamma\)-preserved numerical lattice, and one projects along the rational subspace \(T\). It suffices that the action preserve some full lattice (Li--Zhao Section~2.1); if needed multiply \(H_X\) by a common denominator to put it in that lattice.

For a direct lattice proof, the image of integral Weil divisors in \(V\) is \(\Gamma\)-invariant because
strict transformation permutes prime divisors. To see it is a lattice, take a projective smooth
resolution \(\rho\):Y\(\to X\). Every Weil class is pushed forward from a Cartier divisor on Y. For
numerically trivial \(P\) on Y, the \(\rho\)-exceptional rational divisor \(\rho^*\rho_*P-P\) is
\(\rho\)-numerically trivial, hence zero by the negativity lemma applied to both signs; thus \(\rho_*P\)
is numerically trivial, by projection. The image factors through finitely generated NS(Y). It is
therefore a finitely generated group of rational numerical classes spanning \(V\), hence a full lattice;
\(p\) sends it to a lattice because \(T\) is rational.

The two descriptions of \(\lambda_d\) below prove that a translation with parameter z acts by \(x\mapsto
x+T_{\pi x}z\) with \(T_H=1\). Hence the image translation group in \(G\) is exactly this full lattice,
no larger kernel being counted as an element of the faithful group. A finite extension of this lattice
need not split; no splitting is required.

\subsection*{4. Irregularity one}

Here A is an elliptic curve. For a geometric point \(x\) of \(F\) let \(o_x\):A\(\to F\) be its orbit
map. Pullback takes \(\operatorname{Pic}^0(F)\) to \(A^\vee\) by an isogeny j independent of \(x\): the
induced map on Albanese varieties is the fixed homomorphism A\(\to\)Alb(\(F\)), because the action on the
Albanese torsor is translation by its image. Its dual is j. The integer \(a(d)=\deg(o_x^*M)\) is
independent of \(x\), since the orbit maps form a family over connected \(F\) and degree is constant in a
family of line bundles on A. It is a rational linear functional on \(R\), nonnegative on nef classes and
strictly positive at \(H\) because the orbit map is nonconstant and \(H\) ample.

If \(P_1=\mathcal O_A(0)\), the identity \[
 j\big(t_b^*M-M\big)=t_b^*o_x^*M-o_x^*M
                    =a(d)\lambda_{P_1}(b)
\] follows from \(o_x\circ t_b=t_b\circ o_x\) and the identification of elliptic NS with the degree group. As an equality of homomorphisms on the connected abelian group A it determines \(\lambda_M\) (the finite kernel of j introduces no ambiguity); on \(\mathbb Q\)-vector spaces one can directly invert j. In particular \(\lambda_H\) is an isogeny and on \(U\) one has \[
 T_d=\frac{a(d)}{a(H)}\,I,\qquad [\lambda_d(b)]=T_d[\lambda_H(b)].
\] Use any rational positive definite scalar product on \(U\).

Here is a rational linear section \(s\colon R\to L\) with the correct fibre on all effective faces, including those where the scalar vanishes. Choose \(x\) in \(F\) over a finite Galois extension. Average the pullbacks \(o_{\sigma x}^*\) over its Galois orbit, obtaining a map \(r\colon\operatorname{Pic}(F)_{\mathbb Q}\to\operatorname{Pic}(A)_{\mathbb Q}\) defined over \(k\) (take norms, then divide by the degree). Its restriction on the rational \(Pic^0\) groups is j. Set \[
 w_0(M)=j^{-1}\big(r(M)-a([M])P_1\big).
\] This projection to \(\operatorname{Pic}^0(F)(k)_{\mathbb Q}\) is the identity on \(\operatorname{Pic}^0(F)(k)_{\mathbb Q}\) and determines a rational section \(s\) before and after quotienting by \(K_0\). Suppose \(M\) is effective nef with \(a([M])=0\).

For rational \(M\) some multiple is globally generated by Section 2. Its restriction to every \(o_x\) is
then globally generated and of degree zero, hence trivial. It follows that \(w_0(M)=0\). For real \(M\)
the decomposition in Section 2 expresses it as a nonnegative combination of rational effective nef
classes, each of which has zero orbit degree if their sum does. This proves the same assertion for real
classes and in particular proves that each such \(d\) has the unique effective-nef lift \(s(d)\).

For rational effective nef \(d\) with \(a(d)>0\), take a rational effective nef lift
\(L_d\). Its actual translates are still effective nef, and in the quotient they differ from \(L_d\) by
\(T_d\Lambda\). The convex hull of a full lattice in a real vector space is the entire space, as a
real fundamental parallelepiped and all of its translates are the convex hulls of their vertices. Thus
the entire fibre \(s(d)+U\) occurs. For real \(d\) one can either use its positive
rational-generator decomposition or the general Minkowski argument in Section 6.

\subsection*{5. Abelian surface torsors}

Write \(A^\vee=\operatorname{Pic}^0(F)\); use the ample polarization \(h=\lambda_H:A\to A^\vee\). For
each generic numerical class \(d\) put \(\alpha_d=h^{-1}\lambda_d\in End_k^0(A)\) (extend linearly in
real \(d\)). A rational line bundle on the torsor defines this symmetric homomorphism in the same way as
on A; translation of a chosen geometric origin has no effect on \(\lambda_d\). The symmetry is taken with
respect to the Rosati involution for h. For a geometric affine automorphism \(\varphi\) with linear part
\(u\), pullback on \(Pic^0\) is \(u^\vee\); via h this is \(u^\dagger=h^{-1}u^\vee h\) on the parameter
module. If this operator is in \(G_0\), then \(u^\vee h u=h\) and \(u^\vee\lambda_d u=\lambda_d\) for
each \(d\), so \(u^\dagger=u^{-1}\) and \(u^\dagger\) commutes with \(\alpha_d\). This is consistent with
conjugation of translation pullbacks, which changes \(b\) to \(u^{-1}b\). We always use the ordinary left
action of \(End_k^0(A)\) on \(A(k)_{\mathbb Q}\); dualization itself reverses the order and is not
identified with that left action. To let \(\alpha_d\) act on the visible quotient, identify \(U\) with
the quotient of \(A(k)_{\mathbb Q}\) by \(h^{-1}K_0\) and prove endomorphism invariance as follows.

The class of the torsor \(F\) has finite period \(m\), since it has a point over a finite separable
extension and the restriction-corestriction composite is the degree. For a\(\in End_k\)(A), the
endomorphism \(u\)=1+mn a fixes this torsor class for every integer \(n\).~Take a nonzero \(n\) making
\(u\) an isogeny: the determinant polynomial on a faithful rational Tate representation is nonzero at
zero. Since \(u_*[F]=[F]\), contracting the torsor along \(u\) gives a \(u\)-equivariant \(k\)-morphism
\(F\to F\). It acts as \(u^\vee\) on \(Pic^0\), independently of the translation used to identify
torsors. Extend the generic morphism as an \(S\)-rational map of \(X\). Resolve it as
\(p\)':Z\(\to X\) birational and \(q\)':Z\(\to X\) a morphism. For any rational divisor
\(D\equiv_S0\), \(q'^*D\equiv_S0\).

The divisor \(p'^*p'_*q'^*D-q'^*D\) is \(p\)'-exceptional and \(p\)'-numerically trivial, hence vanishes
by the negativity lemma applied to both signs. Consequently \(p'_*q'^*D\equiv_S0\) (lift any
contracted curve by a generically finite curve on the resolution). Its restriction to \(F\) represents
the generic pullback because the initial generic map is a morphism. Vertical divisors are sent to
vertical divisors in this computation. This proves \(u^\vee\) preserves \(K_0\), and subtraction and
division by mn show \(a^\vee\) preserves it. Duality (or isogeny via h) gives stability of
\(h^{-1}K_0\) under the entire rational endomorphism algebra \(E_A=End_k^0(A)\). In particular
\(\alpha_d\) induces an operator \(T_d\) on \(U\), \(T_H=I\), rational linear in \(d\).~By construction
\([\lambda_d(b)]=T_d[\lambda_H(b)]\) in the identification by \(h\).

For completeness, there is a rational positive definite scalar product on \(U\) making all \(T_d\) for
nef \(d\) self-adjoint and positive semidefinite. The algebra \(E_A\) is semisimple by Poincaré complete
reducibility over \(k\). The finite-dimensional rational \(E_A\)-module \(U\) embeds in the regular left
module \(E_A^n\): a finite generating set gives a surjection and semisimplicity gives a splitting. Choose
on \(E_A\) the form \((v,w)=Tr(vw^\dagger )\), where Tr is the ordinary rational trace on a faithful
homology representation of A after an embedding into \(C\). Equivalently one can take the rational
endomorphism trace of the abelian variety, as in Milne~\cite{Milne}, Abelian Varieties, Thm. 17.3.

This form is positive definite and the adjoint of left multiplication by a is left multiplication by
\(a^\dagger\), by cyclicity of trace. Restrict the sum form on \(E_A^n\) to \(U\). For rational nef
\(d\), on the tangent representation of the complex abelian variety the Hermitian form for \(H\) is
positive definite and the form for \(d\) semipositive (nef classes are limits of ample ones). The
operator \(\alpha_d\) compares the two forms and has only nonnegative real eigenvalues. The exponential
uniformization presents the underlying real tangent space of A(\(C\)) as
H1(A(\(C\)),\(Q\))\(\otimes R\), equivariantly for every endomorphism.

Indeed its period lattice is H1(A(\(C\)),Z), a full lattice in the real tangent space. Hence the
eigenvalues of \(\alpha_d\) on complexified rational homology consist of its complex tangent eigenvalues
and their conjugates, all nonnegative real numbers. Homology is faithful as a rational representation of
\(E_A\), so the rational minimal polynomial on homology annihilates \(\alpha_d\) in \(E_A\) itself and
hence on every rational \(E_A\)-module, including \(U\). Self-adjointness on \(U\) now gives positive
semidefiniteness.

For ample \(d\), the endomorphism is invertible and positive definite. One may embed the finitely
generated field supporting the abelian variety, polarizations, and chosen endomorphism basis into \(C\)
for this argument, or fix a field isomorphism from the algebraic closure of \(\mathbb C(S)\) to \(C\); no
embedding fixing the original copy of \(C\) is required. The real statement on positive combinations of
effective nef classes follows by linearity.

We next supply an initial rational section with the exact effective fibres. In a finite Galois splitting
field choose a point \(p\) of \(F\). Let \(i_p\) be reflection \(x\mapsto2p-x\). On geometric Picard
groups the rational class \(s_p(d)=(M+i_p^*M)/2\) is independent of the lift \(M\) of \(d\), because
\(i_p^*\) acts as -1 on \(Pic^0\) and as +1 on NS. Average over the conjugates of \(p\) and use the norm
to obtain a rational section \(s(d)\) in \(\operatorname{Pic}(F)_{\mathbb Q}\), and then in \(L\).
This construction suffices; it does not require an integral origin or a splitting of the automorphism
group. If \(d\) has an effective rational lift \(M\), the averaged lift \(s(d)\) is rationally
effective: take effective representatives, reflect them after splitting, sum all conjugates and descend
the resulting divisor.

For an effective rational line bundle \(M\) on \(F\) all its translates are nef. Work over the algebraic
closure with any chosen origin and put \(N=\ker(\lambda_d)^0\). The annihilator of \(N\) inside
\(A^\vee\) equals im \(\lambda_d\). Indeed restriction \(i_N^\vee \circ\lambda_d=(\lambda_d\circ
i_N)^\vee =0\) by symmetry, and both connected subvarieties have dimension dim A-dim \(N\) by the dual
exact sequence. In particular translating \(M\) does not change its restriction to \(N\) up to
isomorphism after identifying the coset with \(N\). Such restrictions have numerical class zero, since
the polarization homomorphism of the restriction vanishes. On a general coset a nonzero section does not
vanish identically; it restricts to a nonzero section of a numerically trivial line bundle on the abelian
variety \(N\), which must be trivial.

It follows that \(M\) has trivial restriction on every coset under the translational identifications.
This also holds for its reflected conjugates, hence for the rational section \(s(d)\) on these
cosets. Therefore \(M-s(d)\) is a rational \(\operatorname{Pic}^0\) point annihilating \(N\), i.e.~it lies in im
\(\lambda_d\) on geometric points. Its class lies in \(\lambda_d(A(k)_{\mathbb Q})\): multiply to
remove denominators, take a preimage over a finite extension and take the sum of its Galois conjugates.
This proves that any actual rational effective lift differs from \(s(d)\) by an element of \(\operatorname{im}T_d\)
in \(U\). This implication still applies to real effective nef combinations as in Section 6.

Conversely for rational nef effective \(d\) choose a rational effective lift as above. Since its relative
translations form the coset \(T_d\Lambda\), the convex hull of these effective lifts is
\(s(d)+\operatorname{im}T_d\). \(\operatorname{Conv}(T_d\Lambda)=\operatorname{im}T_d\) follows by applying the linear map \(T_d\)
to \(\operatorname{Conv}(\Lambda)=U\). No discreteness of the image is necessary for this argument.

\subsection*{6. Effective faces and the strong chamber}

In both cases the following hold for every real \(d\) in the generic domain \(B\) and indeed throughout \(D_{\mathrm{gen}}\): \[
 (p(C))\cap\pi^{-1}(d)=s(d)+\operatorname{im}T_d .
\] For the forward inclusion, Section 2 represents any actual lift as a finite nonnegative combination of rational effective nef bundles \(M_i\) whose numerical images \(d_i\) may lie anywhere in \(D_{\mathrm{gen}}\). Their differences from \(s(d_i)\) are in \(\operatorname{im}T_{d_i}\) by Sections 4 and 5 (automatic when the scalar is positive in Section 4). For positive coefficients \(t_i\) the radical of \(\sum_i t_iT_{d_i}\) is the intersection of their radicals, since the forms are positive semidefinite. Taking orthogonal complements gives \(\operatorname{im}T_d=\sum_i\operatorname{im}T_{d_i}\). For the reverse inclusion over a rational generator of \(B\) the translated effective lifts and convexity suffice by the preceding calculations; for real \(d\) in \(B\) take its nonnegative ray-generator expansion and add these full affine fibres. To see the same statement on \(D_{\mathrm{gen}}\) one can instead use any effective rational-nef decomposition. At \(d=0\) the permitted fibre is precisely the zero vector.

We apply the operator-domain lemma proved in the appendix (Lemma~\ref{lem:operator}). Its hypotheses are
verified here:

\begin{enumerate}
\def\labelenumi{\arabic{enumi}.}
\tightlist
\item
  The exact rational sequence in Section 2 has finite dimension, \(G\) acts faithfully on \(L\) through the \(\Gamma\)-image, and it preserves \(p(C)\). It is countable and \(p(C)\) contains the image of the relatively ample open cone. The strong nondegenerate chamber on \(R\) is \(B\).
\item
  There is a rational linear section from Section 4 or 5. The exact real fibre equality holds, including zeros and all boundary faces. The normalized operators are rational linear with \(T_H=I\); they are self-adjoint positive semidefinite on \(B\) in a rational positive definite scalar product. The sum of the ray generators of the full-dimensional \(B\) lies in its interior, hence in the ample cone of \(F\). Its operator is positive definite. Here the ample cone is the interior of the closed nef cone.
\item
  The actual image of \(A(k)\) in \(G_0\) is the normal finite-index subgroup identified with the full lattice \(\Lambda\) by Section 3 and the identities \([\lambda_d(b)]=T_d[\lambda_H(b)]\). The possible finite linear holonomy and affine displacements are handled by Lemma~\ref{lem:operator}. In particular one does not replace \(G_0\) by an arbitrarily selected translation lattice or assume that an extension by the finite point group splits.
\end{enumerate}

To describe the construction and exact separation, average the rational section modulo the translations
\(d\mapsto T_dz\) over the finite quotient \(G_0/\Lambda\) and use an invariant rational
positive scalar product. This yields a section \(s'\) with the same real fibres. The action of
\(g\in G_0\) on a lift \(s'(d)+w\) is \(s'(d)+A_gw+T_dn_g\), where \(A_g\) is
orthogonal, commutes with \(T_d\) and \(n_g\in U_{\mathbb Q}\); the offsets form a finite union \(S\) of
\(\Lambda\)-cosets containing \(\Lambda\).

The closed cone over \(B\) is cut out by \[
 2\langle m,w\rangle\ \le\ \langle m,T_dm\rangle \quad (m\in S).
\] There are only finitely many needed inequalities, by the congruence-class and Dickson-lemma argument in Lemma~\ref{lem:operator} (valid for rational semidefinite ray forms). They force \(w\) perpendicular to ker \(T_d\) on every face because ker \(T_d\) is the rational intersection of the ray radicals in that face. Thus the inequalities add no ineffective boundary directions. Nearest lattice-coset centres in \(U\)/ker \(T_d\) give coverage, including \(d\) on faces. A nonzero offset yields strict separation against the whole translated cone by evaluating the inequality for that offset and the inverse element. A nonidentity base action is separated by \(B\). Only the finite zero-offset subgroup can stabilize the cone.

The finite-group refinement in Section 1 now gives a strong rational polyhedral fundamental domain \(P\)
on \(p(C)\). Section 1 then produces a rational polyhedral domain \(\Pi\subset C\) on \(V\) with
\(\Gamma\Pi=C\), \(\gamma\Pi\cap\)Int \(\Pi\) empty for every nonidentity \(\gamma\), and trivial
setwise stabilizer. Its interior is the full-dimensional relative interior used in Definition 3.4.
\end{proof}

\section*{Further directions}
The proof reduces the strong relative statement to two kinds of input: a
polyhedral chamber for the numerical cone of the generic fibre and a description
of the effective fibres of the Picard quotient. In relative dimension at most two,
surface abundance and the structure of connected automorphism groups supply
these inputs. In larger relative dimension, establishing comparable generic
chambers and controlling the kernels of the induced operators on effective
faces are natural questions. The operator-domain construction is independent of
the surface classification and may apply whenever such a positive-semidefinite
family and a full lattice of actual translations can be identified.

\appendix
\section{An operator domain for a lattice extension}\label{app:operator}
\begin{lemma}\label{lem:operator}
All vector spaces and structural linear maps are defined over \(\mathbb Q\) and are extended to \(\mathbb R\) when discussing cones. Let \[
0\longrightarrow U\longrightarrow L\xrightarrow{\pi} R\longrightarrow0
\] be exact. Let a countable faithful group \(G\subset\mathrm{GL}(L)\) act preserving this sequence and a convex cone \(\mathcal C\subset L\) with nonempty interior. Write \(\bar G\) for the image on \(R\), \(G_0\) for the kernel, and assume \(\pi\mathcal C\) has a strong rational polyhedral domain \(B\) for \(\bar G\), including weak separation against the whole closed translate as in Li--Zhao, Definition 3.4. Let \(T:R\to\operatorname{End}(U)\) be rational linear, \(T_H=1\) at a specified rational point \(H\). Suppose:

\begin{enumerate}
\def\labelenumi{\arabic{enumi}.}
\tightlist
\item
  \(U\) admits a rational positive definite scalar product for which each \(T_D\), \(D\in B\), is self-adjoint and positive semidefinite. Some positive combination of operators at finitely many rational generators of \(B\) is positive definite.
\item
  There is a rational linear section \(s:R\to L\) such that for every \textbf{real} \(D\in B\), \[
  \mathcal C\cap\pi^{-1}(D)=s(D)+\operatorname{im}T_D .
  \]
\item
  A full lattice \(\Lambda\subset U\) is realized as a normal subgroup of finite index in \(G_0\) by \(\tau_z(x)=x+T_{\pi x}z\). The extension of the finite point group by this lattice need not split. No splitting compatible with the group action is assumed in (2).
\end{enumerate}

Then \(\mathcal C\) has a rational polyhedral weak domain with finite stabilizer, and the finite-stabilizer refinement in part 1 of the main proof yields a strong domain. The construction never adjoins directions over faces where \(T_D\) has smaller image.

\end{lemma}
\begin{proof}
\subsection*{Normalization in the kernel of the base action}

For \(g\in G_0\), let \(A_g\) be its restriction to \(U\) and \(B_g\) its conjugation action on the abstract lattice \(\Lambda\), extended linearly to \(U\). Then \(A_g T_D=T_D B_g\) on every \(D\in R\), by applying \(g\tau_z g^{-1}=\tau_{B_gz}\) to a lift of \(D\). At \(H\) this gives \(A_g=B_g\), hence \(A_g\) commutes with every \(T_D\). The \(A_g\) form a finite group since translations act trivially on \(U\). Averaging the original scalar product over this finite group preserves rationality, positivity and self-adjointness: for example each summand \(\langle A_gu,T_D A_gu\rangle\) is nonnegative. We now use the averaged product.

Two rational sections differ by a member of \(\operatorname{Hom}(R,U)\). In that space the subspace \(T(U)=\{D\mapsto T_Du:u\in U\}\) is invariant under the action \(A_g\). The induced affine action of \(G_0\) on the space of rational sections modulo \(T(U)\) factors through \(G_0/\Lambda\). Take representatives \(g_j\) for this finite quotient and form the actual rational section \(s'=\frac1{[G_0:\Lambda]}\sum_j g_j s\). Its class modulo \(T(U)\) is fixed. Accordingly there is a unique \(n_g\in U_\mathbb Q\) such that \[
  g s'(D)=s'(D)+T_D n_g\qquad(D\in R).
\] Uniqueness follows on evaluating at \(H\). Every \(g_js(D)\) lies in \(s(D)+\operatorname{im}T_D\) for \(D\in B\), so \(s'\) still satisfies (2). The group law gives \(n_{gh}=n_g+A_g n_h\), and \(n_{\tau_z}=z\). Therefore \[
S=\{n_g:g\in G_0\}
\] is a finite union of \(\Lambda\)-cosets, contains the zero coset, and is contained in a finite superlattice by rationality. The subgroup \(G_*=\{g\in G_0:n_g=0\}\) embeds into \(G_0/\Lambda\).

\subsection*{Finitely many linear inequalities}

Define \[
\Pi=\bigl\{s'(D)+w:D\in B,\ w\in U,\
 2\langle m,w\rangle\le \langle m,T_Dm\rangle \quad(m\in S)\bigr\}.
\] Write \(B=\sum_i\mathbb R_{\ge0}d_i\) with rational generators and \(q_i(m)=\langle m,T_{d_i}m\rangle\). Every \(q_i\) is rational and positive semidefinite. The sum is positive definite, since a positive linear combination is. Clear denominators for all \(q_i\) on a finite superlattice containing \(S\). Partition \(S\) into its finitely many congruence classes modulo \(2\Lambda\). In each such class, the coordinatewise minimal elements of \(\{(q_i(m))_i\}_m\subset\mathbb N^k\) form a finite set by Dickson's lemma. There are finitely many \(m\) realizing each fixed tuple, since its total \(q=\sum q_i\) has bounded level sets. Retain every such minimal representative.

For any other \(m\), choose \(n\) in the same congruence class with \(q_i(n)\le q_i(m)\) for every \(i\), with one strict. Set \(a=(m+n)/2,\ b=(m-n)/2\). Then \(a\in m+\Lambda\subset S\), \(b\in\Lambda\subset S\), and \[
 a+b=m,\qquad q_i(a)+q_i(b)=\tfrac12(q_i(m)+q_i(n))\le q_i(m).
\] Both \(q(a),q(b)\) are strictly less than \(q(m)\). Induction on the common-denominator integral values of \(q\) proves the inequality for \(m\) from the finite retained inequalities, using \(T_D=\sum t_iT_{d_i}\) for \(D=\sum t_id_i,\ t_i\ge0\). Thus \(\Pi\) is a closed rational polyhedral cone even when some face operators are degenerate.

\subsection*{Exact fibres, covering and stabilizer}

For fixed \(D=\sum t_i d_i\) with \(t_i\ge0\), let \(K_D=\ker T_D\). Positivity gives \(K_D=\bigcap_{t_i>0}\ker T_{d_i}\), so \(K_D\) is rational, including at zero. Since \(S\) contains both signs of every lattice vector of \(K_D\), the inequalities imply \(\langle w,K_D\rangle=0\). Thus \(w\in K_D^\perp=\operatorname{im}T_D\), exactly the permitted fibre directions in (2). In particular the closedness and finite description of \(\Pi\) have not enlarged the effective boundary.

If \(x=s'(D)+T_Du\in\mathcal C\) with \(D\in B\), the inequalities for \(\Pi\) say precisely that \([0]\) in \(U/K_D\) is a nearest centre to \([u]\) among the images of \(S\), using the positive definite squared norm induced by \(\langle u,T_Du\rangle\). As \(K_D\) is rational, the image of \(\Lambda\) is a full lattice of this quotient. The images of \(S\) are a finite union of its cosets; thus a nearest centre \([n_g]\) exists. The affine map \(u\mapsto n_g+A_g u\) descends to a norm isometry on differences in \(U/K_D\) and permutes the centres. Its inverse sends \([u]\) to a point having \([0]\) as nearest centre; equivalently \(g^{-1}x\in\Pi\). For any other base class take an element of \(\bar G\) moving it to \(B\) and lift it to \(G\). This establishes \(G\Pi=\mathcal C\). Countability and nonempty interior of \(\mathcal C\) imply \(\Pi\) is full-dimensional.

For \(g\) with nonidentity image in \(\bar G\), a point of \(g\Pi\cap\operatorname{Int}\Pi\) would project to \((\bar g B)\cap\operatorname{Int}B\). This is empty because \(\pi\) is open and \(B\) satisfies exact strong separation. For \(g\in G_0\) with \(n_g=0\), \(A_g\) preserves \(S\) and all forms, so \(g\Pi=\Pi\). If \(g\in G_0\) with \(m=n_g\ne0\), the functional \[
 s'(D)+w\longmapsto \langle m,T_Dm\rangle-2\langle m,w\rangle
\] is nonzero on \(L\) (vary \(w\)). It is strictly positive on \(\operatorname{Int}\Pi\). For \(x=s'(D)+w\in g\Pi\), apply the inequality for \(n_{g^{-1}}=-A_g^{-1}m\) to \(g^{-1}x=s'(D)+A_g^{-1}(w-T_Dm)\). Orthogonality and commutation give \(2\langle m,w\rangle\ge\langle m,T_Dm\rangle\). Hence \(g\Pi\cap\operatorname{Int}\Pi=\varnothing\). The setwise stabilizer is exactly \(G_*\), which is finite.

\end{proof}

\begin{thebibliography}{99}
\bibitem{LiZhao} Z. Li and H. Zhao, \emph{On the relative Morrison--Kawamata cone conjecture}, Proc. Lond. Math. Soc. \textbf{131} (2025), e70099, doi:10.1112/plms.70099.
\bibitem{LiII} Z. Li, \emph{On the relative Morrison--Kawamata cone conjecture (II)}, J. Algebraic Geom. \textbf{36} (2027), 1--50; theorem numbering here follows arXiv:2309.04673v1 (2023).
\bibitem{HaconXuNumerical} C. D. Hacon and C. Xu, \emph{On finiteness of $B$-representation and semi-log canonical abundance}, arXiv:1107.4149v2.
\bibitem{HaconXuGood} C. D. Hacon and C. Xu, \emph{Existence of log canonical closures}, arXiv:1105.1169v3.
\bibitem{Fujino} O. Fujino, \emph{Minimal model theory for log surfaces}, Publ. RIMS \textbf{48} (2012), 339--371; arXiv:1001.3902v3.
\bibitem{LiFiniteness} Z. Li, \emph{On finiteness of log canonical models}, arXiv:2006.01370.
\bibitem{CKM} H. Clemens, J. Koll\'ar and S. Mori, \emph{Higher-dimensional complex geometry}, Ast\'erisque \textbf{166} (1988).
\bibitem{FujinoGongyo} O. Fujino and Y. Gongyo, \emph{Log pluricanonical representations and the abundance conjecture}, Compos. Math. \textbf{150} (2014), 593--620.
\bibitem{Xu} J. Xu, \emph{Homogeneous fibrations on log Calabi--Yau varieties}, arXiv:1611.01804v2.
\bibitem{Totaro} B. Totaro, \emph{The cone conjecture for Calabi--Yau pairs in dimension two}, Duke Math. J. \textbf{154} (2010), 241--263, doi:10.1215/00127094-2010-039.
\bibitem{Looijenga} E. Looijenga, \emph{Discrete automorphism groups of convex cones of finite type}, Compos. Math. \textbf{150} (2014), 1939--1962, doi:10.1112/S0010437X14007404.
\bibitem{KM} J. Koll\'ar and S. Mori, \emph{Birational geometry of algebraic varieties}, Cambridge Tracts in Math. \textbf{134}, Cambridge University Press, 1998.
\bibitem{Kawamata} Y. Kawamata, \emph{On the cone of divisors of Calabi--Yau fiber spaces}, Internat. J. Math. \textbf{8} (1997), 665--687.
\bibitem{Milne} J. S. Milne, \emph{Abelian varieties}, author's notes, Sections 12 and 17.
\end{thebibliography}
\end{document}
